\section{Introduction}
\label{sec:intro}

Vision--language models (VLMs) are rapidly moving into settings where their inputs are photographs of the real world: assistive agents, content moderation, robotics, and consumer applications. In these settings the model's \emph{safety alignment} --- its ability to recognize that a request is harmful and refuse it --- is a deployment requirement, not an optional extra. Yet safety alignment is almost universally trained and evaluated on \emph{clean} images. Real images are not clean. They arrive blurred, noisy, compressed, poorly lit, and weather-degraded --- the ordinary corruptions catalogued by ImageNet-C~\cite{hendrycks2019benchmarking} that the robustness literature has studied for years, but almost exclusively through the lens of \emph{utility}: does the model still classify, caption, or answer correctly?

This paper asks a different question: \textbf{does a VLM's safety alignment survive ordinary image corruption?} The question matters because in many multimodal safety failures the harm signal lives \emph{in the image}. A prompt like ``how do I make this at home?'' is benign or dangerous depending entirely on what the picture shows. If corruption degrades the safety-relevant visual evidence faster than it degrades general visual competence, then a model can remain perfectly \emph{useful} while silently becoming \emph{unsafe} --- a failure mode that clean-image safety benchmarks cannot detect, and that an attacker can trigger with a Gaussian blur rather than an adversarial perturbation.

We present a systematic study of this question across five open VLMs spanning non-reasoning and reasoning designs --- Llama-3.2-Vision~\cite{llama3herd}, LLaVA-CoT~\cite{xu2024llavacot}, Qwen2.5-VL~\cite{bai2025qwen25vl}, LlamaV-o1~\cite{thawakar2025llamavo1}, and R1-Onevision~\cite{yang2025r1onevision} (with and without its thinking mode) --- under severity-controlled ImageNet-C corruptions, on safety benchmarks chosen to \emph{vary where the harm signal lives}: in image semantics (SIUO~\cite{siuo2024}, HoliSafe~\cite{holisafe2025}, VLSBench~\cite{hu2024vlsbench}), in rendered text (MM-SafetyBench~\cite{liu2024mmsafetybench}, FigStep~\cite{gong2023figstep}), or in the text prompt itself. All conditions share one frozen inference protocol, and safety is scored by a GPT-4o judge that separately rates the harmfulness of the model's \emph{reasoning trace} and of its \emph{final answer}, alongside utility benchmarks (ScienceQA~\cite{lu2022scienceqa}, MathVista~\cite{lu2024mathvista}) measured under the same corruptions.

Our findings are consistent and, we believe, consequential:

\begin{itemize}
  \item \textbf{Ordinary corruptions weaken safety, consistently.} On SIUO, final-answer harmful rate rises in 16 of 18 model$\times$corruption cells (two-sided sign test, $p{=}0.0013$) and in all six configurations under zoom blur ($+3.0$ to $+9.0$ points; the two decreases are statistically indistinguishable from noise). \todo{Replicated across five models on the HoliSafe safe-image+safe-text$\rightarrow$unsafe subset ($n{=}476$): judging in progress.}
  \item \textbf{Safety breaks before capability does.} On the same safety-tuned model, all ten tested corruptions raise attack success ($+1.3$ to $+8.4$, $p{=}0.002$) while ScienceQA accuracy moves at most a few points --- and the corruptions with \emph{zero} utility cost (JPEG, frost, fog) still add up to $+4.2$ points of ASR. The failure is not that the model goes blind; degraded visual evidence simply stops triggering the safety behavior. We term this the \emph{safety--capability decoupling}.
  \item \textbf{The effect is signed, governed by where the harm lives.} Corruption makes models \emph{less} safe when the risk is visible only in the image (VLSBench $+7.0$), but \emph{more} safe when the image itself is the attack vector (typographic MM-SafetyBench $-11.3$): blurring an attack destroys it, while blurring evidence hides it. This \emph{harm-location principle} reconciles apparently contradictory observations about whether corruption helps or hurts VLM safety.
  \item \textbf{Safety tuning cracks under corruption --- and we localize why.} A safety fine-tune that cuts SIUO attack success from 67.7\% to 25.1\% on clean images loses ground under all ten corruptions (worst: 33.5\% under zoom blur) while remaining robust on text-borne harm, exactly as the harm-location principle predicts. A module ablation shows the final refusal decision lives in the LLM, not the vision tower: corruption-aware training of the vision tower yields the safest \emph{reasoning} in our study yet the least safe \emph{answers} --- the reasoning-to-answer pathway is the weakest link.
  \item \textbf{Lightweight mitigations help, unevenly.} Corruption-aware safety fine-tuning heals the reasoning trace (HR $-5.6$) but not the final answer ($+2.4$); a corruption-aware system prompt, by contrast, recovers most of the lost safety at zero training cost (e.g., $74.8\rightarrow 53.9$ under zoom blur), beats a generic safety prompt in every condition, and costs no measurable utility.
\end{itemize}

Together these results identify common image corruption as an overlooked, easily triggered failure mode of multimodal safety alignment --- one invisible to clean-image benchmarks --- and argue that safety evaluation and safety training for deployed VLMs must become corruption-aware.
